\section{Appendix F -- Reviewer Objection Navigation}
% R005_APPENDIX_READER_FRAME
Reader frame. This appendix supports the main manuscript by explaining corpus role, evidence class, audit route, and limitation boundary. It is not a detached raw dump.


\label{sec:oc13-reviewer-navigation}

This appendix replaces the legacy compact route table with prose navigation. The old table was useful as an internal checklist, but in a public monograph it produced broken text extraction and forced the reader to reconstruct sentence order from narrow cells. The publication-grade route below keeps the same reviewer questions while making the reading path explicit.


\subsection{Core theorem route and scope}

A reviewer should begin with the central defended result: bounded post-collapse identity classification with no general lift theorem. The proof-bearing route begins in the results chapter and theorem roadmap, then continues through the proof sheets and finite semantic witnesses. Collapse is formal only when admissible-state emptiness, k-to-zero behavior, and threshold failure are tied to the stated assumptions. Residue is a post-collapse structural trace, not a continuing live identity. Rebirth is a later live continuum grounded in residue and treated as numerically new unless endpoint-bound identity evidence is supplied.


\subsection{Upper-level and hierarchy questions}

The K-level hierarchy is reviewed by asking whether the upper-level roles can be losslessly reduced. If K11 and K12 collapse to K10 without loss of witness, boundary, or reflection behavior, the upper-level claim reopens. If they preserve distinct declared roles, the hierarchy remains a witness taxonomy under the current release scope. This is a reviewer test, not a claim of completed total metaphysics.


\subsection{Journal-core and monograph relation}

The monograph is larger than the journal core because it carries the didactic shell, formal architecture, proof-route explanation, evidence-boundary method, and repair logic. The journal core is smaller because it is an extraction, not a replacement. If the journal core changes theorem scope or hides a required boundary, the extraction fails and must be repaired.


\subsection{External criticism route}

External criticism is treated as pressure on the manuscript, bibliography, evidence rows, replay method, theorem fate, tuple/K rigor, and release package. Acknowledgement of a reviewer records substantive pressure; it never implies endorsement. Predictive claims remain bounded unless the public artifact contains protocol, data or code, output, comparator, and falsifier.


\subsection{Decisive release-failure checks}

A decisive failure is not a vague dislike of style. It is a continuing-identity claim after lawful death, an excluded theorem family quietly entering the proof route, an empirical row without comparator or falsifier being promoted as validation, a same-claim comparator absorbing the residual delta, a broken public PDF, stale metadata, or a public claim that says more than its evidence can carry.
